\chapter{Aspectos Legales y Normativos en Argentina}
\label{sec.cap4:aspectos_legales}

El marco regulatorio argentino en materia de seguridad laboral y protección contra incendios en establecimientos
industriales se asienta en la Ley de Higiene y Seguridad en el Trabajo Nº 19.587, la Ley de Riesgos del Trabajo Nº 24.557
y el Decreto Reglamentario 351/79. En particular, el Capítulo 18 y el Anexo VII de este decreto establecen las
prescripciones de diseño estructural, compartimentación de sectores, medios de escape y cálculo de cargas combustibles.

\section{Decreto Reglamentario 351/79 (Anexo VII)}

La reglamentación exige que todo establecimiento se divida en sectores de incendio aislados entre sí mediante muros
cortafuego y cerramientos cuya resistencia al fuego ($F$) esté garantizada para un lapso temporal determinado (e.g.,
$F30$, $F60$, $F90$, $F120$, $F180$). Esta compartimentación busca confinar el foco inicial, evitando su propagación al
resto del local y posibilitando la evacuación segura del personal antes del colapso estructural.

\subsection{Determinación Analítica de la Carga de Fuego}

Para caracterizar el potencial destructor de un recinto, la normativa define la Carga de Fuego ($Q_f$) como la masa de
madera estándar equivalente por unidad de superficie de piso ($kg/m^2$) capaz de liberar una cantidad de energía térmica
igual a la combustión completa de todos los materiales contenidos en el sector. La fórmula matemática de cálculo pericial
se expresa en la ecuación \ref{eq.cap4:carga_fuego}:

\begin{equation}
\label{eq.cap4:carga_fuego}
Q_f = \frac{\sum (M_i \cdot PC_i)}{S \cdot 4400}
\end{equation}

Donde las magnitudes representan:
\begin{itemize}
    \item $M_i$: Masa total inventariada de cada material combustible presente en el sector ($kg$).
    \item $PC_i$: Poder Calorífico Inferior de cada sustancia considerada ($kcal/kg$).
    \item $S$: Superficie útil de piso del sector de incendio ($m^2$).
    \item $4400$: Constante estandarizada que representa el Poder Calorífico Inferior de la madera de referencia
          ($4400\kcal/\kg \approx 18.4\,\mathrm{MJ/kg}$).
\end{itemize}

El valor de $Q_f$, conjugado con el tipo de riesgo del material dominante (desde Riesgo 1, explosivos, hasta Riesgo 5,
incombustibles), determina la resistencia al fuego mínima de muros y los requerimientos de ventilación y extinción.

\subsection{Unidades de Ancho de Salida (UAS) y Medios de Evacuación}

El dimensionamiento de las vías de escape responde a los criterios deterministas establecidos en el Inciso 3 del Anexo VII
del Decreto Reglamentario 351/79. Una Unidad de Ancho de Salida (UAS) representa el espacio libre requerido para que circule
una hilera de personas en marcha ordenada. Para edificios nuevos, la reglamentación fija que las dos primeras unidades equivalen a
$0.55\m$ de ancho útil cada una (fijando un ancho mínimo inderogable de $2\UAS = 1.10\m$), mientras que cada unidad adicional
subsecuente añade $0.45\m$. Dicho ancho debe ser medido estrictamente entre zócalos, considerándose la cota más restrictiva del
pasaje. Ninguna puerta de escape o pasaje secundario podrá tener un ancho libre inferior a $0.90\m$. En edificios existentes
(previos a abril de 1979), se admiten anchos reducidos transitorios únicamente si resulta técnicamente imposible la ampliación,
conforme al artículo 2 del Anexo I del decreto.

Para calcular la población teórica ($N$) a evacuar, se computa la superficie de piso ($S$) descontando las áreas ocupadas por
los propios medios de escape, sanitarios y locales comunes (para evitar una sobreocupación ficticia del sector), y se divide
por el factor de ocupación ($f_o$), conforme a la ecuación \ref{eq.cap4:ocupantes}:

\begin{equation}
\label{eq.cap4:ocupantes}
N = \frac{S}{f_o}
\end{equation}

El factor $f_o$ ($m^2/\text{persona}$) se determina según el destino edilicio conforme a la tabla del Inciso 3.1.2.
Determinado $N$, el número teórico de unidades de ancho de salida ($n$) se calcula mediante la ecuación \ref{eq.cap4:uas}:

\begin{equation}
\label{eq.cap4:uas}
n = \frac{N}{100}
\end{equation}

Toda fracción igual o superior a $0.5$ se redondea al entero superior más próximo. Así, una salida de $2\UAS$ ($1.10\m$) tiene
capacidad reglamentaria para evacuar entre $1$ y $149$ personas teóricas, y una de $3\UAS$ ($1.55\m$) hasta $349$ personas.

En cuanto a la cantidad de medios de escape independientes ($ME$), la reglamentación estipula que cuando por cálculo corresponda
$n \le 3$, bastará con un único medio de escape o escalera ($ME = 1$). En cambio, cuando corresponda $n \ge 4$, el flujo debe
fraccionarse obligatoriamente en medios de escape y escaleras independientes según la expresión legal \ref{eq.cap4:escape_indep}:

\begin{equation}
\label{eq.cap4:escape_indep}
ME = \frac{n}{4} + 1
\end{equation}

Donde las fracciones iguales o mayores a $0.50$ se redondean a la unidad siguiente. Adicionalmente, rigen las condiciones de
distancia máxima de recorrido: desde el punto más desfavorable de trabajo hasta un medio de escape o salida al exterior no
pueden superarse los $40\m$. En escaleras, los peldaños deben responder a la fórmula ergonómica de paso de Blondel,
establecida en la inecuación \ref{eq.cap4:escalera}:

\begin{equation}
\label{eq.cap4:escalera}
2a + p = 0.60\m \text{ a } 0.63\m
\end{equation}

Donde la alzada ($a$) no puede superar los $0.18\m$ y la pedada ($p$) no debe ser inferior a $0.26\m$.

\section{Reglamentación de la Asociación Electrotécnica Argentina (AEA 90364)}

La reglamentación AEA 90364 establece los lineamientos obligatorios para el proyecto de instalaciones eléctricas en
inmuebles. En materia de seguridad contra incendios destacan dos capítulos:

\begin{itemize}
    \item Sección 42 (Protección contra efectos térmicos): Establece requisitos perentorios sobre canalizaciones y
          conductores en locales con alta densidad de personas o con riesgos de incendio (clasificación BD2, BD3 y BD4).
          Exige el empleo exclusivo de conductores no propagadores de la llama, de reducida emisión de humos opacos y
          libres de halógenos y gases corrosivos (bajo norma IRAM 62267 o IEC 60332), previniendo que la combustión del
          aislamiento propague el siniestro o emita ácido clorhídrico ($HCl$) que asfixie al personal y corroa los
          equipos.
    \item Sección 43 (Protección contra sobrecorrientes): Prescribe la coordinación térmica estricta de las protecciones
          contra cortocircuitos. Para garantizar que un cable no sufra degradación pirolítica durante el tiempo de
          despeje de una falla, se debe satisfacer la inecuación de energía pasante \ref{eq.cap4:tiempo_corte}:
\end{itemize}

\begin{equation}
\label{eq.cap4:tiempo_corte}
I_{cc}^2 \cdot t \le k^2 \cdot S^2
\end{equation}

Donde $I_{cc}$ representa la corriente de cortocircuito eficaz presunta en el punto de falla en amperes, $t$ el tiempo
de actuación de la protección en segundos, $S$ la sección nominal útil del conductor en $mm^2$, y $k$ una constante
térmica dependiente del metal (cobre o aluminio) y del tipo de aislamiento (e.g., $k = 115$ para cobre con aislación de
PVC y $k = 143$ para cobre con aislación de polietileno reticulado XLPE). Si esta condición no se verifica, el cable
alcanza su temperatura límite de régimen de cortocircuito ($160\Celsius$ para PVC), iniciando la carbonización y la
generación de llamas.

\section{Ejemplos Prácticos de Aplicación y Cálculo Reglamentario}
\label{sec.cap4:ejemplos_practicos}

Con el fin de ilustrar la aplicación concreta de las prescripciones legales analizadas, se desarrollan a continuación
tres memorias de cálculo representativas de la práctica pericial en instalaciones industriales y de ingeniería
electrónica.

\subsection{Ejemplo 1: Memoria de Cálculo de Carga de Fuego y Dotación de Extintores}
\label{sec.cap4:ejemplo_carga_fuego}

Se considera un depósito anexo a una nave de montaje de equipos electrónicos cuya superficie útil de piso es $S = 600\m^2$.
El relevamiento de inventario de materiales combustibles arroja las siguientes masas y poderes caloríficos inferiores:

\begin{itemize}
    \item Embalajes celulósicos (cajas de cartón corrugado): masa $M_1 = 3200\kg$, con poder calorífico
          $PC_1 = 4000\kcal/\kg$.
    \item Tarimas de soporte logístico (pallets de pino): masa $M_2 = 2000\kg$, con poder calorífico
          $PC_2 = 4400\kcal/\kg$.
    \item Plásticos de embalaje (polietileno expandido): masa $M_3 = 600\kg$, con poder calorífico
          $PC_3 = 9500\kcal/\kg$.
    \item Gabinetes poliméricos (resinas ABS y policarbonato): masa $M_4 = 800\kg$, con poder calorífico
          $PC_4 = 8000\kcal/\kg$.
\end{itemize}

El calor total acumulado en el sector se calcula integrando la energía liberable por cada material:

\begin{align}
\sum (M_i \cdot PC_i) &= (3200 \cdot 4000) + (2000 \cdot 4400) + (600 \cdot 9500) + (800 \cdot 8000) \notag \\
&= 12800000 + 8800000 + 5700000 + 6400000 = 33700000\kcal
\end{align}

Aplicando la ecuación \ref{eq.cap4:carga_fuego}, se obtiene la carga de fuego ponderada equivalente en madera:

\begin{equation}
Q_f = \frac{33700000\kcal}{600\m^2 \cdot 4400\kcal/\kg} = \frac{33700000}{2640000} \approx 12.77\kg/\m^2
\end{equation}

A partir de este resultado analítico, se contrastan las exigencias reglamentarias del Anexo VII del Decreto 351/79:

\begin{enumerate}
    \item Clasificación de Riesgo: El sector corresponde a Riesgo 3 (Muy Combustible), dado el predominio de plásticos
          y material celulósico de embalaje.
    \item Resistencia al Fuego de Muros: Conforme al Cuadro 2.2.1 del Anexo VII, para un establecimiento de Riesgo 3 con
          ventilación natural y carga de fuego $Q_f \le 15\kg/\m^2$, los muros perimetrales de compartimentación deben
          poseer una resistencia al fuego mínima $F60$ (sesenta minutos de resistencia mecánica y estanqueidad a humos).
    \item Dotación de Extintores Portátiles: De acuerdo al Cuadro 4 del Anexo VII, para Riesgo 3 y $Q_f \le 15\kg/\m^2$,
          se exige una unidad extintora de potencial mínimo $2\mathrm{A}$ por cada $200\m^2$ de superficie de piso. El número
          mínimo de extintores requeridos resulta:
          \begin{equation}
          N_{\text{ext}} = \frac{600\m^2}{200\m^2} = 3 \text{ extintores Clase A}
          \end{equation}
          Debido a la presencia de tableros eléctricos seccionales y luminarias energizadas, se añade la exigencia de
          capacidad de extinción para riesgos Clase C. En consecuencia, se instalan tres extintores de polvo químico seco
          ABC (fosfato monoamónico) de $5\kg$ con rating certificado $4\mathrm{A}:40\mathrm{B:C}$, superando ampliamente el
          potencial normativo. Asimismo, se distribuyen de modo tal que la distancia máxima de recorrido desde cualquier
          punto operativo hasta el matafuego más cercano no exceda los $20\m$.
\end{enumerate}

\subsection{Ejemplo 2: Dimensionamiento de Medios de Escape y Salidas de Emergencia}
\label{sec.cap4:ejemplo_uas}

Se proyecta una nave de manufactura electrónica de dos plantas con una superficie útil de $S = 1600\m^2$ en planta alta,
destinada a líneas de inserción superficial (SMT) y bancos de control de calidad.

\begin{enumerate}
    \item Cálculo de Población Teórica ($N$): Para establecimientos industriales, el Decreto 351/79 estipula un factor de
          ocupación $f_o = 16\m^2/\text{persona}$. Mediante la ecuación \ref{eq.cap4:ocupantes}:
          \begin{equation}
          N = \frac{1600\m^2}{16\m^2/\text{persona}} = 100 \text{ personas}
          \end{equation}
    \item Determinación de Unidades de Ancho de Salida ($n$): Aplicando la ecuación \ref{eq.cap4:uas}:
          \begin{equation}
          n = \frac{100}{100} = 1\UAS
          \end{equation}
          Sin embargo, el marco legal impone la restricción perentoria de que ninguna salida colectiva puede tener un ancho
          inferior a $2\UAS$. Por lo tanto, se adoptan obligatoriamente $2\UAS$, correspondientes a un ancho libre de paso
          de $1.10\m$.
    \item Supuesto de Alta Densidad y Salidas Múltiples: Si la nave albergase un sector de capacitación anexo o un área de
          empaque manual con una población recalculada de $N = 320$ operarios, las unidades de ancho de salida resultarían:
          \begin{equation}
          n = \frac{320}{100} = 3.2 \implies 4\UAS \quad (\text{redondeo al entero superior})
          \end{equation}
          El ancho útil total demandado se calcula sumando las dos primeras unidades ($1.10\m$) y las dos adicionales
          ($2 \times 0.45\m = 0.90\m$), alcanzando $2.00\m$ de luz libre. Al verificarse que $n \ge 4$, el flujo de personas
          no puede evacuarse por una sola puerta, rigiendo la exigencia de medios de escape independientes según la
          ecuación \ref{eq.cap4:escape_indep}:
          \begin{equation}
          ME = \frac{4}{4} + 1 = 2 \text{ medios de escape independientes}
          \end{equation}
          Esto impone la apertura de dos salidas diferenciadas de $2\UAS$ ($1.10\m$) cada una, ubicadas en extremos opuestos
          de la planta para impedir que un foco ígneo anule simultáneamente ambas vías de evacuación, verificando que la
          distancia máxima hacia cualquiera de ellas no supere los $40\m$.
    \item Verificación Ergonómica de la Escalera de Evacuación: La caja de escaleras presurizada que conecta con la planta
          baja se diseña con peldaños de alzada $a = 0.17\m$ y pedada $p = 0.28\m$. Verificando la relación de Blondel
          (\ref{eq.cap4:escalera}):
          \begin{equation}
          2a + p = 2(0.17\m) + 0.28\m = 0.34\m + 0.28\m = 0.62\m
          \end{equation}
          El valor obtenido se sitúa dentro del intervalo reglamentario $[0.60\m; 0.63\m]$, garantizando un descenso fluido
          y seguro del personal en condiciones de emergencia.
\end{enumerate}

\subsection{Ejemplo 3: Verificación Térmica de Conductores ante Cortocircuito}
\label{sec.cap4:ejemplo_cortocircuito}

En un tablero general de baja tensión ($400\,\mathrm{V}$ trifásico, $50\,\mathrm{Hz}$), se proyecta un alimentador hacia un
centro de distribución secundario que abastece bancos de potencia y servidores. El alimentador está constituido por cables
unipolares de cobre con aislamiento termoplástico de PVC según norma IRAM NM 247-3, con una sección nominal de fase
$S = 25\mm^2$.

El estudio de cortocircuito arroja en las barras de salida una corriente de cortocircuito trifásica simétrica presunta
$I_{cc} = 8500\,\mathrm{A}$ ($8.5\,\mathrm{kA}$). El circuito está protegido aguas arriba por un interruptor automático en
caja moldeada (MCCB) con unidad de disparo magnético regulada, cuyo tiempo total de apertura y extinción del arco eléctrico
ante fallas francas es $t = 30\,\mathrm{ms} = 0.03\,\mathrm{s}$.

A fin de comprobar si el conductor se encuentra debidamente protegido contra incendios originados por calentamiento
adiabático, se evalúa la inecuación de energía específica pasante (\ref{eq.cap4:tiempo_corte}):

\begin{enumerate}
    \item Energía Pasante del Cortocircuito ($I_{cc}^2 \cdot t$):
          \begin{equation}
          I_{cc}^2 \cdot t = (8500\,\mathrm{A})^2 \cdot 0.03\,\mathrm{s} = 72.25 \times 10^6 \cdot 0.03 = 2.1675 \times 10^6\,\mathrm{A^2 s}
          \end{equation}
    \item Capacidad Térmica Límite del Conductor ($k^2 \cdot S^2$):
          Para conductores de cobre aislados en PVC, la constante térmica establecida por la reglamentación AEA 90364 es
          $k = 115\,\mathrm{A\sqrt{s}/mm^2}$ (correspondiente a una temperatura de régimen permanente de $70\Celsius$ y un
          límite de cortocircuito de $160\Celsius$). Reemplazando para la sección de $25\mm^2$:
          \begin{equation}
          k^2 \cdot S^2 = (115)^2 \cdot (25\mm^2)^2 = 13225 \cdot 625 = 8.2656 \times 10^6\,\mathrm{A^2 s}
          \end{equation}
    \item Evaluación de la Condición Térmica:
          \begin{equation}
          2.1675 \times 10^6\,\mathrm{A^2 s} \le 8.2656 \times 10^6\,\mathrm{A^2 s}
          \end{equation}
          La energía solicitada por la falla representa únicamente el $26.2\percent$ de la capacidad admisible del cable. Se
          comprueba analíticamente que la aislación polimérica no alcanzará su umbral de degradación pirolítica, impidiendo
          el cebado de llamas sobre las bandejas portacables.
    \item Consecuencia ante Falta de Selectividad o Falla de la Protección:
          Si la protección magnética no operase por descalibración o trabamiento mecánico, obligando a que la sobrecorriente
          sea despejada por el bimetal térmico en un tiempo retardado de $t = 250\,\mathrm{ms} = 0.25\,\mathrm{s}$, la
          energía disipada ascendería a:
          \begin{equation}
          I_{cc}^2 \cdot t = (8500\,\mathrm{A})^2 \cdot 0.25\,\mathrm{s} = 18.0625 \times 10^6\,\mathrm{A^2 s} > 8.2656 \times 10^6\,\mathrm{A^2 s}
          \end{equation}
          En este escenario anómalo, la energía liberada excede en más de un $218\percent$ el umbral admisible, provocando la
          fusión instantánea de la cubierta de PVC, la emisión masiva de gas clorhídrico tóxico y la ignición abierta del
          tendido eléctrico, evidenciando el papel crítico que desempeña el cálculo de protecciones según la reglamentación
          AEA en la prevención de siniestros industriales.
\end{enumerate}

